\documentclass[titlepage,11pt]{article}
\usepackage{comment}
\usepackage{enumitem}
\usepackage{transparent} % Untuk transparansi gambar
\usepackage{listings}
\usepackage{amsmath}
\usepackage{graphicx}
\usepackage[font=small,labelfont=bf]{caption}
\usepackage[bahasa]{babel}
\usepackage{float}
\usepackage{verbatim}
\usepackage{graphicx,tabularx,multirow}
\usepackage{xcolor}
\usepackage[onehalfspacing]{setspace}
\usepackage[
	allcolors=visigrey,
	colorlinks=true,
]{hyperref}
\sloppy
\usepackage[a4paper,left=2cm,right=2cm]{geometry}
% Pengaturan kutipan artikel
\usepackage[style=ieee, backend=biber]{biblatex}
%Code listing style pak akok
\usepackage{listings}

\definecolor{codegreen}{rgb}{0,0.6,0}
\definecolor{codegray}{rgb}{0.5,0.5,0.5}
\definecolor{codepurple}{rgb}{0.58,0,0.82}
\definecolor{backcolour}{rgb}{0.95,0.95,0.92}

\usepackage{eso-pic} % Untuk menambahkan elemen ke seluruh halaman

\newcommand\BackgroundPic{
  \put(0,0){
    \parbox[b][\paperheight]{\paperwidth}{
      \vfill
      \centering
      \transparent{0.1}
      \includegraphics[width=0.4\paperwidth,keepaspectratio]{miot.png}
      \vfill
    }
  }
}

\newcommand\BackgroundAllPages{ \AddToShipoutPicture*{\BackgroundPic} }
\newcommand\BackgroundNone{ \ClearShipoutPicture } % hilangkan background

\lstdefinestyle{mystyle}{
	backgroundcolor=\color{backcolour}, commentstyle=\color{codegreen},
	keywordstyle=\color{magenta},
	numberstyle=\small\color{codegray},
	stringstyle=\color{codepurple},
	basicstyle=\ttfamily\footnotesize,
	breakatwhitespace=false,         
	breaklines=true,                 
	captionpos=t,                    
	keepspaces=true,                 
	numbers=left,                    
	numbersep=5pt,                  
	showspaces=false,                
	showstringspaces=false,
	showtabs=false,           
	frame = single,
	tabsize=2
}
\lstset{style=mystyle}

\definecolor{visigrey}{rgb}{.1,.15,.15}
\geometry{top=1cm,bottom=.5cm}
\savegeometry{titlepage}
\geometry{top=2cm,bottom=2cm}
\savegeometry{main}

\def\bspace{\(\qquad\qquad\qquad\)}
\usepackage[T1]{fontenc}
\usepackage[utf8]{inputenc}
\usepackage{tgheros}
\renewcommand*\familydefault{\sfdefault}

\setcounter{tocdepth}{6}

\def\autor{Laboratorium }
\def\lab{Multimedia dan Internet of Things}
\def\departemen{Departemen Teknik Komputer}
\def\institut{Institut Teknologi Sepuluh Nopember}
\def\praktikum{Laporan Sementara \\ Praktikum Jaringan Komputer}
\def\nama{Ahmad Faiq Fawwaz - 5024231032}
% Ubah Judul sesuai dengan modul
\def\judul{tunnelling, Ipsec, dan Queue}
\def\tanggal{2025}
\begin{document}
% Ubah Bahasa sesuai dengan keinginan
\selectlanguage{bahasa}

\BackgroundNone
\def\headingtype{\bf \small}
\loadgeometry{titlepage}

\begin{titlepage}
	\centering
	\begin{tabularx}{\textwidth}{l@{\hskip 0pt}lX}
		\raisebox{-0.5\height}{\includegraphics[width=3cm]{Cover/img/logodepart.png}} 
		& \raisebox{-0.5\height}{\includegraphics[width=3cm]{Cover/img/miot.png}} 
		& \raggedleft
	\hfill
	\begin{minipage}{0.5\textwidth}
		\raggedleft
		{\emph{\headingtype \autor}} \\[-2pt]
		{\headingtype \lab} \\[-2pt]
		{\headingtype \departemen} \\[-2pt]
		{\headingtype \emph{\institut}}
	\end{minipage}

	\vspace{5cm}
	\end{tabularx}
	
	\vspace{5cm}
	{\Huge \bf \praktikum \par}
	
	\vspace{2cm}
	{\LARGE \bf \expandafter\string\judul \par}
	
	\vspace{2cm}
	{\Large \nama \par}
	
	\vfill
	{\Large \tanggal \par}
	
	\vfill
	\includegraphics[width=\textwidth]{Cover/img/footer.png}
\end{titlepage}

\loadgeometry{main}


\BackgroundAllPages
% Pilih Modul yang akan di build
\section{Langkah-Langkah Percobaan}
Bagian ini menjelaskan secara rinci tahapan atau prosedur yang dilakukan selama praktikum. Langkah-langkah ditulis secara urut dan sistematis, mulai dari persiapan alat hingga pelaksanaan percobaan. Penulisan harus jelas agar dapat dipahami oleh orang lain yang membaca laporan ini.

\section{Analisis Hasil Percobaan}
Berisi pembahasan terhadap hasil yang diperoleh selama praktikum. Analisis dilakukan dengan membandingkan hasil percobaan terhadap teori, serta mengidentifikasi faktor-faktor yang memengaruhi hasil, seperti kesalahan alat atau langkah percobaan. Tujuan bagian ini adalah untuk mengevaluasi keberhasilan percobaan dan memastikan pemahaman praktikan terkait praktikum yang telah dilaksanakan.

\section{Hasil Tugas Modul}
Bagian ini memuat hasil dari pengerjaan tugas tambahan yang diberikan dalam modul praktikum serta penjelasan dari jawaban tersebut. 

\section{Kesimpulan}
Kesimpulan berisi ringkasan dari hasil praktikum dan hal-hal penting yang didapatkan. Bagian ini menjawab tujuan praktikum, mencantumkan hasil yang sesuai atau tidak sesuai dengan teori, serta pembelajaran yang diperoleh oleh praktikan.

\section{Lampiran}
\subsection{Dokumentasi saat praktikum}
Menampilkan foto selama pelaksanaan praktikum. Dokumentasi meliputi foto alat yng digunakan dan foto praktikan saat praktikum. Tujuannya sebagai bukti telah dilakukan kegiatan praktikum.

\subsection{Hasil Challenge Modul}
Memuat hasil dari challenge modul saat praktikum (jika berhasil mengerjakan challenge) serta penjelasan lengkap dari hasil tersebut. 
\textcolor{red}{Jika kelompok praktikum tidak berhasil mengerjakan challenge modul, maka sub section ini dapat dihapus.}


\end{document}